\chapter{热力学第一定律}
\problem{4.1}{试在 $p$-$V$ 图上画出以下理想气体系统完成的准静态过程：
（1）$p=kV$；（2）$p=kT$；（3）$V=kT$，其中 $k$ 为满足各自要求的常量。并计算当系统体积由 $V_1$ 变至 $V_2$ 时，上述三过程中系统对外界所作的功。}
\answerarea
\problem{4.2}{$0.020\,\mathrm{kg}$ 的氮气的温度由 $17\,^\circ\mathrm C$ 上升到 $27\,^\circ\mathrm C$，若在升温过程中：（1）体积保持不变；（2）压强保持不变；（3）不与外界交换热量。试在将这些氮气近似为理想气体的情况下，分别求出气体内能的改变、吸收的热量、外界对气体所作的功。}
\answerarea
\clearpage
\problem{4.3}{一气体系统的定压摩尔热容 $C_p$ 随温度改变的规律为
\[
C_p=a+bT-cT^2,
\]
其中 $a,b,c$ 是常量，当 $\dfrac m\mu\,\mathrm{mol}$ 气体经一个等压过程使温度从 $T_1$ 变到 $T_2$ 时，气体从外界吸收的热量为多少？}
\answerarea
\problem{4.4}{在标准状况下，$1\,\mathrm{mol}$ 的单原子分子理想气体先经过一绝热过程，再经过一等温过程，最后压强和体积均增为原来的两倍，试确定整个过程中气体吸收的热量 $Q_1$。若先经过等温过程再经过绝热过程而达到同样的状态，结果是否相同？若不同，吸收的热量 $Q_2$ 是多少？}
\answerarea
\clearpage
\problem{4.5}{设 $1\,\mathrm{mol}$ 固体的状态方程可以表示为 $V=V_0+aT+bp$；内能可以表示为 $U=cT-apT$，其中 $a,b,c$ 和 $V_0$ 均是常量。试确定：
\begin{enumerate}
\item 摩尔焓的表达式；
\item 摩尔热容量 $C_{p,\mathrm m}$ 和 $C_{V,\mathrm m}$。
\end{enumerate}}
\answerarea
\problem{4.6}{铅弹射击到木板上，经与木板相撞后熔化掉。已知撞击前铅弹的温度是 $30\,^\circ\mathrm C$，而且已知铅的熔点是 $327\,^\circ\mathrm C$，铅的熔解热是 $2.045\times10^4\,\mathrm{J/kg}$，比热是 $130\,\mathrm{J/(kg\cdot K)}$，试确定此铅弹撞击木板前的最低速度是多少？}
\answerarea
\clearpage
\problem{4.7}{$1.0\,\mathrm g$ 氮气原来的温度和压强分别为 $423\,\mathrm K$ 和 $5.066\times10^5\,\mathrm{Pa}$，经准静态绝热膨胀后，体积变为原来的两倍，试确定在这一过程中气体对外界作的功。}
\answerarea
\problem{4.8}{对密度为 $10^4\,\mathrm{kg/m^3}$、质量为 $0.10\,\mathrm{kg}$ 的一块金属在保持温度恒定情况下加压，使压强由 0 逐渐增至 $100\,\mathrm{atm}$，如果该过程中金属的等温压缩系数保持为常量 $6.72\times10^{-12}/\mathrm{Pa}$，试确定该过程中外界对金属所作的功。}
\answerarea
\clearpage
\problem{4.9}{记理想气体的状态方程为 $p=K\overline\varepsilon$，其中 $p$ 为系统的压强，$\overline\varepsilon$ 为平均动能密度，$K$ 为常量，试导出系统的绝热压缩系数。}
\answerarea
\problem{4.10}{在原子弹爆炸后 $0.1\,\mathrm s$ 所出现的“火球”是半径约 $15\,\mathrm m$、温度为 $3000\,\mathrm K$ 的气体球，试确定温度变为 $300\,\mathrm K$ 时，“火球”的半径大约为多大？}
\answerarea
\clearpage
\problem{4.11}{在宇宙大爆炸理论中，起初局限于小区域的辐射能量以球对称方式绝热膨胀，在膨胀过程中逐渐冷却。仅从热力学角度，试推导出温度 $T$ 和辐射半径 $R$ 间的关系。}
\answerarea
\problem{4.12}{两端开口的 U 形管中注入水银，直到水银柱全长为 $h$ 的状态。
\begin{enumerate}
\item 若将一边管中水银压下，然后使水银振荡，试证明：不计摩擦时，振动周期 $T_1$ 为 $T_1=2\pi\sqrt{\dfrac h{2g}}$；
\item 把管的左端封闭起来，使被封在管内的气柱的高度为 $l$，然后使水银柱振荡，假设摩擦可忽略，空气可以近似为理想气体，且可以认为气体经历的过程是准静态绝热过程，而气压计中水银柱的高度为 $h_0$，试证明这种情况下的周期为
\[
T_2=2\pi\sqrt{\frac h{2g+\gamma h_0g/l}}；
\]
\item 试证明绝热指数 $\displaystyle\gamma=\frac{2l}{h_0}\left[\left(\frac{T_1}{T_2}\right)^2-1\right]$。
\end{enumerate}}
\answerarea
\clearpage
\problem{4.13}{记大气的平均摩尔质量为 $\overline\mu$，定压摩尔热容量为 $C_{p,\mathrm m}$，地面 $h=0$ 处的压强、温度分别为 $p_0,T_0$，试证明：按绝热大气模型，大气高度 $h$ 与压强 $p$ 的关系为
\[
h=\frac{C_{p,\mathrm m}T_0}{\overline\mu g}
\left[1-\left(\frac p{p_0}\right)^{(\gamma-1)/\gamma}\right].
\]}
\answerarea
\problem{4.14}{观测表明，距地面不太远的大气的温度随高度而降低，但在称为对流层顶的“边界”，在对流层顶以外，不再具有上述规律。
\begin{enumerate}
\item 试通过具体计算，说明通过测量对流层顶以下一定高度处的气温可以确定该处大气的稳定性。
\item 某日某地区的大气含有饱和水蒸气，地面气温为 $27\,^\circ\mathrm C$，$5000\,\mathrm m$ 高空处气温为 $-5.0\,^\circ\mathrm C$，假设大气可以近似为理想气体，水的摩尔汽化热近似为常量 $4.3\times10^4\,\mathrm J$，饱和大气中水蒸气的摩尔浓度随温度的变化率近似为常量 $0.00021/\mathrm K$，试通过计算说明该地区该日的大气是否稳定。
\end{enumerate}}
\answerarea
\clearpage
\problem{4.15}{地球上有很多雪山，常年积雪不化的最低海拔高度称为雪线高度。假设常温下蒸气的凝结热近似为常量 $2.4\times10^6\,\mathrm{J/kg}$，大气中水蒸气的摩尔浓度随温度的变化率近似为常量 $2.1\times10^{-4}/\mathrm K$，海平面上大气的温度为 $27\,^\circ\mathrm C$，试确定雪线高度。}
\answerarea
\clearpage
\problem{4.16}{如图所示，在一山峰的两侧建有气象观测站 $M_1,M_4$，山顶建有观测站 $M_3$，一侧山坡上建有观测站 $M_2$。通常，潮湿空气可以沿着山坡持续上升并流到另一侧。某日，观测站 $M_1$ 测得其附近的大气压强为 $100\,\mathrm{kPa}$、气温为 $20\,^\circ\mathrm C$，$M_4$ 测得大气压强也为 $100\,\mathrm{kPa}$，$M_3$ 测得的大气压强为 $70\,\mathrm{kPa}$，$M_2$ 测得大气压强为 $84.5\,\mathrm{kPa}$，并发现在其所处高度开始有云形成。潮湿空气继续上升，经约 25 分钟后到达山顶，并且在这段过程中其中的水蒸气凝结成雨而落下。记每平方米面积上空的潮湿空气的质量为 $2000\,\mathrm{kg}$，每 kg 潮湿空气凝结出 $2.45\,\mathrm g$ 雨水，水的汽化热为 $2.5\,\mathrm{kJ/kg}$，并且空气可以近似为理想气体。
\begin{enumerate}
\item 假设空气密度随高度增加而线性减小，试确定云底距离地面的高度。
\item 试确定云底处的气温。
\item 试确定山顶的观测站 $M_3$ 测得的气温。
\item 假设在观测站 $M_2$ 与 $M_3$ 之间的降雨均匀，并记每小时内每平方米面积内降落 $1\,\mathrm{kg}$ 雨水为 $1\,\mathrm{mm}$ 的降雨量，试确定这两个观测站之间的降雨量。
\item 在 $M_4$ 处进行观测，测得的气温很可能为多少 $^\circ\mathrm C$？与 $M_1$ 处的相比，这里的空气有何特点？
\end{enumerate}
\fig{4_16}}
\answerarea
\clearpage
\problem{4.17}{试证明：在泊松比 $\gamma$ 为常数的情况下，如果理想气体在某一过程中的热容量也是常量，则这个过程一定是多方过程。}
\answerarea
\problem{4.18}{一定量的氢气，压强为 $p_1=1.0\,\mathrm{atm}$、体积为 $V_1=2.3\,\mathrm L$、温度为 $t_1=26\,^\circ\mathrm C$；经过一个多方过程，达到压强 $p_2=0.50\,\mathrm{atm}$、体积为 $V_2=4.1\,\mathrm L$。试确定：
\begin{enumerate}
\item 多方指数 $n$；
\item 内能的变化量 $\Delta U$；
\item 对外界所做的功 $W$；
\item 吸收的热量 $Q$。
\end{enumerate}}
\answerarea
\clearpage
\problem{4.19}{证明若理想气体按 $V=a_0/\sqrt p$ 的规律膨胀，其中 $a_0$ 为常量，则气体在该过程中的热容 $C$ 可以表示为
\[
C=C_V-\frac{a_0^2}{TV}.
\]
并说明气体的温度随膨胀而变化的行为。}
\answerarea
\problem{4.20}{引力理论表明，引力塌缩时会辐射出热能。假设一星体的密度为常量，总质量为 $M$，
\begin{enumerate}
\item 试确定其半径由 $R_1$ 减小到 $R_2$ 的过程中释放出的热量；
\item 假定该过程满足过程方程 $V=a_0/\sqrt p$，并且组成星体的物质仍可近似为理想气体，试确定该星体的比热容。
\end{enumerate}}
\answerarea
\clearpage
\problem{4.21}{一圆筒装有压强为 $2\,\mathrm{atm}$ 的氧气，其容积为 $3\,\mathrm L$，温度为 $300\,\mathrm K$，使氧气依次经历下列过程：等压下加热到 $500\,\mathrm K$，等体下冷却到 $250\,\mathrm K$，等压下冷却到 $150\,\mathrm K$，等体下加热到 $300\,\mathrm K$。
\begin{enumerate}
\item 在 $p$-$V$ 图上画出上述四个过程，并给出每一过程的终态的 $p,V$ 值；
\item 计算氧气在一个循环中所作的净功；
\item 计算此循环的效率。
\end{enumerate}}
\answerarea
\problem{4.22}{设有一卡诺热机，它的低温热源的温度为 $280\,\mathrm K$，效率为 $40\%$，现若使此热机效率提高到 $50\%$，试确定：
\begin{enumerate}
\item 如果低温热源的温度保持不变，高温热源的温度必须提高多少？
\item 如果高温热源的温度保持不变，那么低温热源的温度必须降低多少？
\end{enumerate}}
\answerarea
\clearpage
\problem{4.23}{一个平均输出功率为 $50\,\mathrm{MW}$ 的发电厂，热机循环的高温热源温度为 $T_1=1000\,\mathrm K$，低温热源温度为 $T_2=300\,\mathrm K$，试确定：
\begin{enumerate}
\item 理论上，该热机的最高效率为多少？
\item 这个发电厂的热效率实际只能达到理论最大值的 $70\%$，为了产生 $50\,\mathrm{MW}$ 的电功率，每秒钟需要提供多少焦耳的热量？
\item 如果低温热源由一条河流来担当，其流量为 $10\,\mathrm{m^3/s}$，试求由于电厂释放的热量而引起的温升是多少？
\end{enumerate}}
\answerarea
\problem{4.24}{$1\,\mathrm{mol}$ 空气（可视为理想气体）进行如图所示的循环，其中 $C$—$A$ 为绝热过程，试确定：
\begin{enumerate}
\item 循环功；
\item 循环效率；
\item 以循环中最高温度及最低温度为热源的可逆卡诺循环的效率。
\end{enumerate}
\fig{4_24}}
\answerarea
\clearpage
\problem{4.25}{设有一以理想气体为工作物质的热机循环，如图所示，由等压过程 $CA$、等体过程 $AB$ 和绝热过程 $BC$ 组成。若已知 $p_A,p_B,V_A,V_C$ 和绝热指数 $\gamma$，试证明此热机的循环效率为
\[
\eta=1-\gamma\frac{(V_C/V_A)-1}{(V_C/V_A)^\gamma-1}.
\]
\fig{4_25}}
\answerarea
\problem{4.26}{如图，燃料电池是把化学能直接转为电能的装置，把氢气和氧气连续通入多孔 Ni 电极，Ni 电极浸在 KOH 电解液中。在两极进行的化学反应为：
\[
\begin{aligned}
\text{正极：}\quad&2\mathrm{H_2}+4\mathrm{OH^-}\longrightarrow4\mathrm{H_2O}+4e^-；\\
\text{负极：}\quad&2\mathrm{H_2O}+4e^-+\mathrm{O_2}\longrightarrow4\mathrm{OH^-}.
\end{aligned}
\]
液态水在 $25\,^\circ\mathrm C$ 情况下的标准生成焓为 $-285.84\,\mathrm{kJ/mol}$，在此状态下电池的电动势为 $1.229\,\mathrm V$，试求此燃料电池的效率。
\fig{4_26}}
\answerarea
\clearpage
\problem{4.27}{一制冷机的工作物质进行如图所示的循环过程，其中 $AB,CD$ 分别是温度为 $T_2,T_1$ 的等温过程；$BC,DA$ 为等压过程。设工作物质为理想气体，试证明此制冷机的制冷系数为
\[
\varepsilon=\frac{T_1}{T_2-T_1}.
\]
\fig{4_27}}
\answerarea
\problem{4.28}{某理想卡诺制冷机工作在 $0\,^\circ\mathrm C$ 与 $100\,^\circ\mathrm C$ 的水之间，已知冰的熔解热为 $334\,\mathrm{kJ/kg}$，汽化热为 $2260\,\mathrm{kJ/kg}$，试确定要产生 $1\,\mathrm{kg}$ 的冰，需将多少 $100\,^\circ\mathrm C$ 的水变为蒸汽？}
\answerarea
\clearpage
\problem{4.29}{四冲程汽油机的工作循环（奥托循环或定体加热循环）如图 4.21（a）所示。工作开始时，一定量的雾状汽油和空气的混合气体由状态 $V_1,T_A$ 绝热压缩到状态 $V_2,T_B$，达到可燃点；此刻，火花塞放出电火花，点燃气体，经等体加热至状态 $V_2,T_C$；燃烧形成的高温气体经绝热膨胀至状态 $V_1,T_D$；后经等体放热使气体恢复到初始状态 $V_1,T_A$。若工作气体可近似为理想气体，试确定该工作循环的效率。\bookfig{4_29}{图4.21（a）\quad 奥托循环}}
\answerarea
\problem{4.30}{四冲程柴油机的工作循环（狄塞尔循环或定压加热循环）如图 4.22（a）所示。工作开始时，一定量的空气由状态 $V_1,T_A$ 经绝热压缩至状态 $V_2,T_B$；此刻将雾状柴油喷入气缸，燃油与空气混合后燃烧，等压吸热至状态 $V_3,T_C$；然后绝热膨胀至状态 $V_1,T_D$；再经等体放热至状态 $V_1,T_A$。若工作物质可近似为理想气体，试确定该工作循环的效率。\bookfig{4_30}{图4.22（a）\quad 狄塞尔循环}}
\answerarea
\clearpage
\problem{4.31}{记光子气体满足的斯特藩—玻尔兹曼定律中的常量为 $\sigma$，光子数密度 $n$ 与温度间的函数关系 $n=n(T)$ 中的常量为 $a$。
\begin{enumerate}
\item 记光子气体的初始平衡态的状态参量温度和体积分别为 $T_0,V_0$，试导出在准静态绝热过程中光子气体的温度与体积 $V$ 间的函数关系，以及在温度分别为 $T_0,T$（假设 $T<T_0$）的恒温热源之间以光子气体为工作物质的卡诺热机的效率。
\item 如果光子气体由初始状态 $\{T_0,V_0\}$ 出发，经一个自由膨胀，在极短的时间内形成一个体积为 $8V_0$ 的瞬时非平衡态，又在绝热等体条件下，经一段时间后达到体积为 $8V_0$ 的平衡态，记瞬时非平衡态的光子数是初始时的 $\eta_1$ 倍、达到的平衡态的光子数是初始时的 $\eta_2$ 倍，试确定 $\eta_1$ 和 $\eta_2$。
\item 如果光子气体由初始状态 $\{T_0,V_0\}$ 出发经过一个 $\mathrm dQ=\beta\,\mathrm dU$ 的准静态过程达到体积为 $8V_0$ 的平衡态，试确定该过程中光子气体吸收的热量和达到的温度。
\end{enumerate}}
\answerarea
\problem{4.32}{设服从克劳修斯方程 $p(V_0-b)=RT$ 的气体的摩尔内能为 $U_{\rm m}=C_{V,\rm m}T+U_0$，其中 $C_{V,\rm m}$ 和 $U_0$ 都是常量。
\begin{enumerate}
\item 试证明在准静态绝热过程中气体遵从绝热方程 $p(V_0-b)^\gamma=\text{常量}$，其中 $\gamma=C_{p,\rm m}/C_{V,\rm m}$。
\item 试证明利用克劳修斯气体所作的卡诺循环的效率为 $\eta=1-T_2/T_1$。
\end{enumerate}}
\answerarea
\clearpage
\problem{4.33}{某制冷机在冬天作热泵用，给某房间供暖，在室外为 $-10\,^\circ\mathrm C$ 情况下，使房间内的温度保持为 $20\,^\circ\mathrm C$。
\begin{enumerate}
\item 若该房间向外散热主要由面向室外的面积为 $5\,\mathrm{m^2}$、厚度为 $2\,\mathrm{mm}$ 的玻璃引起，设该玻璃的热导率为 $0.75\,\mathrm{W/(m\cdot K)}$，电费收费标准为每度 0.5 元，试计算该空调一天工作 12 小时情况下至少需要花多少电费。
\item 若将上述玻璃换为每层厚度仍为 $2\,\mathrm{mm}$、其中间有 $0.5\,\mathrm{mm}$ 空气夹层的“双层玻璃”，假设空气的热导率为 $0.025\,\mathrm{W/(m\cdot K)}$，电费收费标准不变，试给出该空调机仍然一天工作 12 小时情况下可以节省的电费。
\end{enumerate}}
\answerarea
\clearpage
